\section{Elementary geometry}
\label{sec:geometry}

The first bound relates the number of stones in a window to the maximum gap
from an empty cell to those stones.

\begin{lemma}[Packing bound]
\label{lem:L1}
In a window with $k\ge1$ stones of one colour and no stones of the other,
every empty cell of the window is within distance $6-k$ of a same-window
stone.  The bound is tight.
\end{lemma}

\begin{proof}
Index the cells $0,\ldots,5$ along the window axis.  Let $e$ be empty, and let
$\delta$ be its axis-distance from a nearest in-window stone.  The $\delta$
cells from $e$ toward that stone, including $e$ and excluding the stone, are
empty and lie in the window.  There are $6-k$ empty cells in total, so
$\delta\le6-k$.  By \cref{fact:F1}, axis-distance agrees with hex distance
for these cells.  For tightness, put the stones at offsets
$6-k,\ldots,5$ and take $e$ at offset $0$.
\end{proof}

Figure~\ref{fig:packing} displays the extremal arrangement and a concrete
$k=3$ instance of the general formula.

\begin{figure}[tbp]
  \centering
  \begin{tikzpicture}[scale=.68]
    \HexPatch{0/0,1/0,2/0,3/0,4/0,5/0}
    \HexWindow[window one]{0}{0}{1}{0}
    \HexAttacker{3}{0}\HexAttacker{4}{0}\HexAttacker{5}{0}
    \node[inner sep=0pt] at (2.598,1.52)
      {$\overbrace{\hspace{3.53cm}}$};
    \node[font=\scriptsize] at (2.598,2.08) {$d=6-k=3$};
    \node[font=\scriptsize] at (0,-1.35) {$e$};
    \foreach \q in {0,...,5}{
      \node[font=\scriptsize,yshift=-13mm] at \HexAt{\q}{0} {$\q$};
    }
  \end{tikzpicture}
  \caption{The tightness configuration for the packing bound, shown at
  $k=3$.  In general the $k$ stones occupy offsets $6-k$ through $5$, while
  the empty cell at offset $0$ is at distance $6-k$ from the nearest stone.}
  \label{fig:packing}
\end{figure}

A mixed set of occupied and empty cells on a finite line segment must change
state across some adjacent pair.

\begin{lemma}[Boundary pair]
\label{lem:L2}
Any window containing at least one stone and at least one empty cell has an
empty cell at distance $1$ from a same-window stone.
\end{lemma}

\begin{proof}
The six cells are consecutive on a line.  The stone cells form a nonempty
proper subset of this sequence.  At a boundary of that subset, a stone and an
empty cell are consecutive.  Their hex distance is $1$ by
\cref{fact:F1}.
\end{proof}

The packing bound places all tactics that can resolve within the current turn
near stones already on the board.

\begin{theorem}[Single-turn tactical locality]
\label{thm:T1}
Every placement that completes a win within the current turn, and every empty
cell of every threat window, lies within distance $2$ of an existing stone.
\end{theorem}

\begin{proof}
Completion within at most two placements requires a window with at least four
stones of the mover already present.  Apply \cref{lem:L1} with $k=4$ to its
empty cells.  A threat window also has at least four stones by
\cref{def:D5}, so the same bound applies to every one of its empty cells.
\end{proof}

Creating a threat can start one cell farther away because the placement need
only raise an existing three-stone window to count four.

\begin{theorem}[Threat-creation locality]
\label{thm:T2}
A placement creating an $X$-threat lies within distance $3$ of an $X$ stone.
The bound is tight.
\end{theorem}

\begin{proof}
Immediately before the placement, the window that becomes an $X$-threat held
at least three $X$ stones.  Apply \cref{lem:L1} with $k=3$ to the placement
cell.  For tightness, place $X$ stones at offsets $3,4,5$ of a window and make
the threat-creating placement at offset $0$.
\end{proof}
